\chapter{Conclusión}
\label{cap:conclusion}

El desarrollo del compilador TinyS fue un buen proyecto para poner en
práctica lo aprendido en varias materias de la carrera, como Ingeniería
de Software, Paradigmas de Programación, Lenguajes Formales y
Computabilidad. Durante el trabajo fue necesario organizar el código,
elegir cómo resolver cada parte y preparar pruebas para comprobar que
todo funcionara. También permitió ver cómo se relacionan temas que se
habían estudiado por separado.

El proyecto ayudó a entender mejor las bases de la programación y el
porqué de muchas cosas que usamos al escribir código. Por ejemplo, cómo
se guardan los valores en memoria, qué ocurre cuando se llama a un
método y cómo se detectan los errores de un programa. Tener que resolver
estas cuestiones dentro del compilador hizo que fueran más fáciles de
comprender.

Construir un compilador para TinyS permitió seguir el camino desde un archivo de código
fuente hasta las instrucciones que ejecuta la máquina. Fue una forma de
darle sentido a lo aprendido durante la carrera y de entender mejor qué
hay detrás de los programas que escribimos.
